% !TeX program = lualatex
% ============================================================================
% presentation.tex — Apresentação de Defesa de TCC (EMC/CTC/UFSC)
% Compilação: cd documents_tcc/tex/thesis/presentation && latexmk presentation.tex
% ============================================================================
\documentclass[aspectratio=169,10pt]{beamer}

% Guarda de compilador: utiliza LuaLaTeX (mesma stack de second_draft.tex).
\usepackage{iftex}
\ifLuaTeX\else
  \PackageError{presentation}{Esta apresentacao deve ser compilada com LuaLaTeX}{Use: latexmk presentation.tex}
\fi

\usepackage{ufsc-metropolis}
\usepackage{ufsc-components}

% Metadados institucionais (sincronizados com second_draft.tex)
\title[Paralelização em GPU de AG+2-opt para o TSP]{%
  Avaliação de Estratégias de Paralelização em GPU para um Algoritmo Memético Aplicado ao Problema do Caixeiro Viajante%
}
\tipotrabalho{Defesa de Trabalho de Conclusão de Curso — Engenharia Mecânica (EMC/UFSC)}
\author[Lucas E. G. Silva]{Lucas Eduardo Galdino Silva}
\orientador{Prof.~Dr.~Sérgio Fernando Mayerle}
\institute[UFSC / EMC]{Universidade Federal de Santa Catarina — Centro Tecnológico — EMC}
% Substitua pela data exata da defesa quando agendada (ex.: Florianópolis, 10 de dezembro de 2026)
\date{Florianópolis, 2026}

\hypersetup{
  pdftitle={Avaliação de Estratégias de Paralelização em GPU para um Algoritmo Memético Aplicado ao Problema do Caixeiro Viajante},
  pdfauthor={Lucas Eduardo Galdino Silva},
  pdflang={pt-BR}
}

\begin{document}

% --- Capa institucional UFSC/EMC (alinhamento superior seguro) ---------------
\begin{frame}[plain,noframenumbering,t]
  \titlepage
\end{frame}

% --- Roteiro / Sumário (recomendado pelo regulamento EMC-UFSC) --------------
\begin{frame}{Roteiro da Apresentação}
  \begin{columns}[T,onlytextwidth]
    \begin{column}{0.60\textwidth}
      \tableofcontents[hideallsubsections]
    \end{column}
    \begin{column}{0.37\textwidth}
      \begin{block}{Estrutura da Defesa}
        \begin{itemize}
          \item \textbf{Exposição:} 20--25 min\\{\scriptsize(limite EMC-UFSC: 30 min)}
          \item \textbf{Escopo:} \totalvariantes\ variantes, \totalinstancias\ instâncias TSPLIB, \totalexecucoes\ execuções
          \item \textbf{Arguição:} slides de apoio técnico após as Referências
        \end{itemize}
      \end{block}
    \end{column}
  \end{columns}
\end{frame}

% --- Capítulos da apresentação (1 arquivo por capítulo da monografia) -------
% ============================================================================
% slides/01-introducao.tex — Capítulo 1: Introdução e Objetivos
% ============================================================================

\section{Introdução e Objetivos}

% --- Contexto e Motivação ---------------------------------------------------
\begin{frame}{Contexto e Motivação}
  \begin{columns}[T,onlytextwidth]
    \begin{column}{0.52\textwidth}
      \begin{block}{Problema do Caixeiro Viajante (TSP)}
        \begin{itemize}
          \item Encontrar o ciclo hamiltoniano de menor custo total.
          \item Problema \textbf{NP-difícil} central em otimização combinatória \autocite{cook2012pursuit}.
          \item Base controlada (toda permutação é viável; ótimos conhecidos na TSPLIB).
        \end{itemize}
      \end{block}
    \end{column}
    \begin{column}{0.45\textwidth}
      \begin{block}{Aplicações na Engenharia Mecânica}
        \begin{itemize}
          \item Sequenciamento de furação em chapas e placas (PCBs).
          \item Trajetórias de usinagem CNC e corte a laser \autocite{cook2012pursuit}.
          \item Programação de robôs de solda, inspeção e logística fabril.
        \end{itemize}
      \end{block}
    \end{column}
  \end{columns}

  \vspace{0.18cm}
  \begin{exampleblock}{Oportunidade com GPUs (CUDA) e Desafio de Desempenho}
    \begin{itemize}
      \item Algoritmos meméticos (\textbf{AG + $2$-opt}) combinam busca global e refinamento local, mas têm alto custo para $n$ crescente \autocite{larranaga1999genetic,croes1958method}.
      \item GPUs reúnem centenas/milhares de ALUs e alta largura de banda para paralelismo em dados \autocite{nvidia2024cuda}, porém lançamentos de \emph{kernel} e tráfego PCIe podem anular o ganho \autocite{schulz2013gpu,van2013gpu}.
    \end{itemize}
  \end{exampleblock}
\end{frame}

% --- Problema de Pesquisa e Justificativa -----------------------------------
\begin{frame}{Problema de Pesquisa e Justificativa}
  \begin{itemize}
    \item Estudos comparativos de CPU vs.\ GPU costumam alterar \textbf{operadores, vizinhanças ou critérios de parada} entre plataformas \autocite{schulz2013gpu,van2013gpu,benaini2018genetic}, impedindo isolar o efeito da arquitetura.
  \end{itemize}

  \vspace{0.12cm}
  \begin{alertblock}{Questão Central de Pesquisa}
    \emph{Como avaliar, sob a mesma especificação algorítmica e boas práticas de implementação em cada plataforma, o impacto de diferentes estratégias de paralelização em GPU sobre o desempenho de um algoritmo memético (AG+$2$-opt) para o TSP?}
  \end{alertblock}

  \vspace{0.12cm}
  \begin{columns}[T,onlytextwidth]
    \begin{column}{0.48\textwidth}
      \begin{block}{Paridade de Especificação}
        Mesmos operadores, hiperparâmetros, lista $k_{\text{viz}} = \kviz$, teto $\kappa = \maxpassadas$ e parada — variantes \textbf{análogas}, respeitando boas práticas de CPU (NumPy) e GPU (memória compartilhada).
      \end{block}
    \end{column}
    \begin{column}{0.49\textwidth}
      \begin{block}{Rigor Estatístico}
        Instâncias padronizadas da \textbf{TSPLIB} \autocite{reinelt1991tsplib}, \totalrepeticoes\ repetições pareadas por semente e protocolo não paramétrico de \textcite{demsar2006statistical}.
      \end{block}
    \end{column}
  \end{columns}
\end{frame}

% --- Objetivos --------------------------------------------------------------
\begin{frame}{Objetivos}
  \begin{block}{Objetivo Geral}
    Quantificar o impacto de diferentes estratégias de paralelização em GPU sobre o \textbf{tempo}, o \textbf{\emph{speedup}} e a \textbf{qualidade (\emph{gap})} de um algoritmo memético (AG+$2$-opt) para o TSP, situando-as na taxonomia de \textcite{crainic2003parallel,crainic2010parallel}.
  \end{block}

  \vspace{0.12cm}
  {\small\textbf{Objetivos Específicos:}}
  \begin{enumerate}
    \small
    \item Implementar \textbf{\totalvariantes\ variantes análogas} (\emph{CPU}, \emph{Híbrida Ingênua}, \emph{Híbrida Otimizada}, \emph{Totalmente em GPU}) sob o mesmo algoritmo e critérios de parada.
    \item Avaliar \textbf{\totalinstancias\ instâncias da TSPLIB} (\menorinstancia\ a \maiorinstancia\ cidades; formatos \texttt{EUC\_2D}, \texttt{ATT}, \texttt{GEO} e \texttt{EXPLICIT}).
    \item Aplicar protocolo experimental reprodutível (\totalrepeticoes\ repetições, sementes \sementeinicial--\sementefinal) e testes não paramétricos \autocite{demsar2006statistical}.
    \item Analisar compromissos práticos entre \textbf{tempo, qualidade, tráfego PCIe e VRAM} em uma GPU de entrada.
  \end{enumerate}
\end{frame}


% ============================================================================
% slides/02-fundamentacao.tex — Capítulo 2: Fundamentação Teórica
% ============================================================================

\section{Fundamentação Teórica}

% --- Algoritmo Memético (AG + 2-opt) ----------------------------------------
\begin{frame}{Algoritmo Memético (AG + $2$-opt)}
  \begin{columns}[c,onlytextwidth]
    \begin{column}{0.53\textwidth}
      \begin{itemize}
        \small
        \item \textbf{Exploração global (AG):} $n_{\text{pop}} = 256$ rotas; seleção por torneio ($\tau = 5$), cruzamento OX, mutação \emph{swap} e elitismo $(\mu + \lambda)$ \autocite{goldberg1989genetic,larranaga1999genetic}.
        \item \textbf{Busca local $2$-opt (\emph{best-improvement}):} remove $(v_i,v_{i+1})$ e $(v_j,v_{j+1})$, reconecta e reverte o segmento intermediário \autocite{croes1958method}:
              \[
                \Delta C = (c_{i,j} + c_{i+1,j+1}) - (c_{i,i+1} + c_{j,j+1}) < 0
              \]
        \item \textbf{Lista estática de candidatos ($k_{\text{viz}} = \kviz$):} reduz a varredura de $\Theta(n^2)$ para $\Theta(n\,k_{\text{viz}})$ sem alterar a função objetivo \autocite{lin1973efficient}, limitada a $\kappa = \maxpassadas$ varreduras.
      \end{itemize}
    \end{column}
    \begin{column}{0.44\textwidth}
      \slideTikz[\slideTallFigH]{%
        \begin{tikzpicture}[node distance=3.4mm]
          \node[sl_term, text width=4.6cm] (init) {Inicialização: $P(0)$ ($n_{\text{pop}}=256$)};
          \node[sl_cpu, below=of init, text width=4.6cm] (gen) {Seleção $\to$ Cruzamento OX $\to$ Mutação \emph{Swap}};
          \node[sl_mem, below=of gen, text width=4.6cm] (ls) {\textbf{Busca Local $2$-opt}\\($k_{\text{viz}}=\kviz$, até $\kappa=\maxpassadas$ varreduras)};
          \node[sl_cpu, below=of ls, text width=4.6cm] (eval) {Avaliação + Elitismo $(\mu+\lambda)$};
          \node[sl_dec, below=4mm of eval, text width=2.3cm] (stop) {Parada atingida?};
          \node[sl_good, below=4mm of stop, text width=3.6cm] (end) {Retorna melhor circuito};

          \draw[sl_arr] (init) -- (gen);
          \draw[sl_arr] (gen) -- (ls);
          \draw[sl_arr] (ls) -- (eval);
          \draw[sl_arr] (eval) -- (stop);
          \draw[sl_arr] (stop) -- node[sl_lbl, right=2pt] {Sim} (end);
          \draw[sl_arr] (stop.east) -- ++(5.5mm,0) |- node[sl_lbl, pos=0.25, right=1pt] {Não ($t{+}1$)} (gen.east);
        \end{tikzpicture}%
      }{Elaborado pelo autor.}
    \end{column}
  \end{columns}
\end{frame}

% --- Contraste Arquitetural CPU vs. GPU -------------------------------------
\begin{frame}{Contraste Arquitetural: CPU vs.\ GPU \autocite{flynn1972taxonomy}}
  \begin{columns}[T,onlytextwidth]
    \begin{column}{0.48\textwidth}
      \slideFigure[\slideDualFigH]{\linewidth}{cpu-gpu-silicon-allocation.png}{Adaptado de \textcite{paz2011gpucpuimage}.}
      \vspace{0.1cm}
      \begin{block}{CPU (MIMD / SIMD)}
        Poucos núcleos complexos, amplo controle e \emph{caches} grandes — otimizada para \textbf{baixa latência} em fluxos sequenciais.
      \end{block}
    \end{column}
    \begin{column}{0.49\textwidth}
      \slideFigure[\slideDualFigH]{\linewidth}{gpu-internal-architecture.png}{Adaptado de \textcite{shah2023gpuarchitectureimage}.}
      \vspace{0.1cm}
      \begin{block}{GPU (SIMT / CUDA)}
        Centenas de ALUs compactas e memória global de alta vazão — otimizada para \textbf{\emph{throughput}} paralelo em dados.
      \end{block}
    \end{column}
  \end{columns}
\end{frame}

% --- Modelo SIMT e Lei de Amdahl --------------------------------------------
\begin{frame}{Modelo SIMT e Gargalo Sequencial (Lei de Amdahl)}
  \begin{columns}[T,onlytextwidth]
    \begin{column}{0.52\textwidth}
      \begin{block}{Execução em GPU (Modelo SIMT) \autocite{nvidia2024cuda}}
        \begin{itemize}
          \item \emph{Threads} escalonadas em \textbf{\emph{warps} de $32$ unidades} executando a mesma instrução.
          \item \textbf{Divergência de fluxo:} desvios condicionais (\texttt{if/else}) irregulares serializam o \emph{warp}.
          \item \textbf{Coalescimento:} leituras contíguas em memória global evitam múltiplas transações.
          \item \textbf{Memória compartilhada:} reduz latência dentro de cada SM (\emph{Streaming Multiprocessor}).
        \end{itemize}
      \end{block}
    \end{column}
    \begin{column}{0.45\textwidth}
      \begin{alertblock}{Lei de Amdahl \autocite{amdahl1967validity}}
        \[
          S = \frac{1}{(1 - f) + \dfrac{f}{p}} < \frac{1}{1 - f}
        \]
        \begin{itemize}
          \item $f$: fração paralelizável; $p$: aceleração obtida.
          \item Ao acelerar intensamente apenas a busca local ($p \gg 1$), a fração restante $(1 - f)$ em CPU passa a dominar o tempo total.
        \end{itemize}
      \end{alertblock}
    \end{column}
  \end{columns}
\end{frame}


% ============================================================================
% slides/03-metodologia.tex — Capítulo 3: Materiais e Métodos
% ============================================================================

\section{Materiais e Métodos}

% --- Ambiente Computacional (Hardware e Software) ---------------------------
\begin{frame}{Ambiente Computacional}
  \begin{slidetabela}
    \begin{tabular}{@{}lll@{}}
      \toprule
      \textbf{Recurso} & \textbf{CPU (Host)} & \textbf{GPU (Device)} \\
      \midrule
      Processador / Placa & Intel Core i7-7700HQ        & NVIDIA GeForce GTX 1050 Mobile \\
      Arquitetura         & Kaby Lake (2017)            & Pascal GP107 (Compute Capability 6.1) \\
      Unidades            & 4 núcleos / 8 \emph{threads} & 640 núcleos CUDA em 5 SMs \\
      Frequência          & 2,80--3,80 GHz              & 1354--1493 MHz \\
      Memória             & 16 GB RAM DDR4              & 4 GB VRAM GDDR5 ($\approx 112$ GB/s) \\
      Barramento CPU--GPU & \multicolumn{2}{c}{PCIe 3.0 ($\approx 16$ GB/s)} \\
      \bottomrule
    \end{tabular}
  \end{slidetabela}
  \notatabelaSlide{Elaborado pelo autor (Seção~3.1.1 da monografia).}

  \vspace{0.18cm}
  \begin{block}{\emph{Stack} de Software e \emph{Kernels} Customizados}
    \begin{itemize}
      \item \textbf{Plataforma:} Debian GNU/Linux 13, Python 3.12.13, NumPy 2.5.4 (CPU) e CuPy 14.2.0 / CUDA 12.9 (GPU).
      \item \textbf{Totalmente em GPU:} usa \textbf{\emph{kernels} CUDA C++ customizados} (\texttt{RawKernel}) pré-compilados para fundir etapas, usar memória compartilhada e cooperar entre \emph{threads}.
    \end{itemize}
  \end{block}
\end{frame}

% --- Restrições Físicas e Aritmética Inteira --------------------------------
\begin{frame}{Fronteiras Físicas e Aritmética Inteira Exata}
  \begin{columns}[T,onlytextwidth]
    \begin{column}{0.48\textwidth}
      \begin{block}{Gargalo do Barramento PCIe ($\sim 7\times$)}
        \begin{itemize}
          \item Largura de banda interna da VRAM ($\approx 112$ GB/s) é cerca de \textbf{7 vezes maior} que a do PCIe ($\approx 16$ GB/s):
                \[
                  T_{\text{total}} = T_{\text{cpu}} + T_{\text{gpu}} + T_{\text{transf}}
                \]
          \item Transferências frequentes penalizam diretamente o tempo total.
        \end{itemize}
      \end{block}
    \end{column}
    \begin{column}{0.49\textwidth}
      \begin{block}{Memória VRAM e Precisão Exata}
        \begin{itemize}
          \item Em \texttt{pr1002} ($n=\maiorinstancia$), todo o estado residente ocupa apenas \textbf{$6{,}07$ MB} dos 4 GB de VRAM.
          \item Rotas até $n=2048$ cabem inteiras na memória compartilhada dos SMs.
          \item \textbf{Aritmética inteira:} $D$ em \texttt{int32} ($\text{nint}$) e custos em \texttt{int64} eliminam desvios de ponto flutuante \autocite{reinelt1991tsplib}.
        \end{itemize}
      \end{block}
    \end{column}
  \end{columns}
\end{frame}

% --- Arquitetura Modular e Hiperparâmetros Fixos ----------------------------
\begin{frame}{Arquitetura Modular e Hiperparâmetros Fixos}
  \begin{columns}[c,onlytextwidth]
    \begin{column}{0.45\textwidth}
      \slideTikz[\slideTallFigH]{%
        \begin{tikzpicture}[node distance=4mm]
          \node[sl_gpu, text width=5.2cm] (c1) {\textbf{Camada 1 — Algoritmos}\\\texttt{GeneticAlgorithm} + 6 estratégias intercambiáveis (CPU/GPU) $\to$ \totalvariantes\ variantes};
          \node[sl_cpu, below=5mm of c1, text width=5.2cm] (c2) {\textbf{Camada 2 — Orquestração}\\\texttt{TSPLIB}, pré-cálculo ($D$, $N_k$), \totalrepeticoes\ sementes, \emph{checkpoints} JSON $\to$ DuckDB};
          \node[sl_good, below=5mm of c2, text width=5.2cm] (c3) {\textbf{Camada 3 — Análise Estatística}\\Protocolo Demšar, tabelas \LaTeX, diagramas $CD$ e gráficos PDF};
          \draw[sl_arr] (c1) -- node[sl_lbl, right=2pt] {Execução controlada} (c2);
          \draw[sl_arr] (c2) -- node[sl_lbl, right=2pt] {Telemetria por fase} (c3);
        \end{tikzpicture}%
      }{Elaborado pelo autor.}
    \end{column}
    \begin{column}{0.52\textwidth}
      \begin{slidetabeladensa}[\slideTallFigH]
        \begin{tabular}{@{}ll@{}}
          \toprule
          \textbf{Parâmetro} & \textbf{Valor fixo (4 variantes)} \\
          \midrule
          População ($n_{\text{pop}}$)       & $256$ rotas ($< 1{,}03$ MB até $n=\maiorinstancia$) \\
          Seleção ($\tau_{\text{trnmt}}$)    & Torneio de tamanho $5$ \\
          Cruzamento / Mutação               & OX ($p_c = 1{,}0$) / \emph{Swap} ($p_{\text{mut}} = 0{,}02$) \\
          Busca local $2$-opt                & \emph{Best-improvement}, $\kappa \le \maxpassadas$ varreduras \\
          Lista de candidatos                & $k_{\text{viz}} = \kviz$ vizinhos mais próximos \\
          Sobrevivência                      & Elitismo $(\mu + \lambda)$, melhores $256$ de $512$ \\
          Parada (i) alvo                    & $\text{gap} \le \varepsilon_{\text{alvo}} = 1\%$ frente a $B^*$ \\
          Parada (ii) estagnação             & $G_{\text{estag}} = \lceil 2\sqrt{n} \rceil$ gerações sem melhora \\
          Parada (iii) teto                  & $G_{\max} = \lceil 2\,n\sqrt{n} \rceil$ gerações \\
          \bottomrule
        \end{tabular}
      \end{slidetabeladensa}
      \notatabelaSlide{$k_{\text{viz}}=\kviz$ acelera a busca em até $8{,}22\times$ (CPU) e $4{,}05\times$ (GPU) (ver Backup~2).}
    \end{column}
  \end{columns}
\end{frame}

% --- As Quatro Variantes de Paralelização -----------------------------------
\begin{frame}{As Quatro Variantes de Paralelização (Tipo 1)}
  \slideTikz[\slideTikzMaxH]{%
    \begin{tikzpicture}
      \node[font=\footnotesize\bfseries\color{ufscDark}] at (0, 0) {Variante};
      \node[font=\footnotesize\bfseries\color{ufscSlate}] at (3.6, 0) {CPU (Host)};
      \node[font=\footnotesize\bfseries\color{ufscAccent}] at (6.7, 0) {Barramento PCIe};
      \node[font=\footnotesize\bfseries\color{ufscBlue}] at (10.3, 0) {GPU (VRAM / CUDA)};

      \node[sl_term, text width=2.2cm] (v1) at (0, -0.95) {\textbf{1. CPU}\\(\emph{baseline})};
      \node[sl_cpu, text width=3.5cm] (c1) at (3.6, -0.95) {Todo o ciclo AG + $2$-opt\\(NumPy vetorizado / SIMD)};
      \node[font=\scriptsize\color{ufscSlate}] at (6.7, -0.95) {---};
      \node[sl_cpu, dashed, text width=3.9cm, text=ufscSlate] (g1) at (10.3, -0.95) {Não utilizada};

      \node[sl_term, text width=2.2cm] (v2) at (0, -2.15) {\textbf{2. Híbrida}\\\textbf{Ingênua}};
      \node[sl_cpu, text width=3.5cm] (c2) at (3.6, -2.15) {Sel., OX, \emph{Swap}, Aval.,\\Elitismo $(\mu+\lambda)$};
      \node[sl_mem, text width=3.9cm] (g2) at (10.3, -2.15) {$2$-opt (\textbf{1 rota / \emph{kernel}})};
      \draw[sl_arr_pcie] (c2.east) -- node[above=1pt, font=\tiny\bfseries\color{ufscAccent}] {256$\times$ rota + $D$ / ger.} node[below=1pt, font=\tiny\color{ufscSlate}] {H2D + D2H repetitivo} (g2.west);

      \node[sl_term, text width=2.2cm] (v3) at (0, -3.35) {\textbf{3. Híbrida}\\\textbf{Otimizada}};
      \node[sl_cpu, text width=3.5cm] (c3) at (3.6, -3.35) {Sel., OX, \emph{Swap},\\Elitismo $(\mu+\lambda)$};
      \node[sl_gpu, text width=3.9cm] (g3) at (10.3, -3.35) {$2$-opt + Avaliação em \textbf{lote}\\($D$ e $N_k$ fixas em VRAM)};
      \draw[sl_arr_pcie] (c3.east) -- node[above=1pt, font=\tiny\bfseries\color{ufscBlue}] {1 matriz $P(t)$ / ger.} node[below=1pt, font=\tiny\color{ufscSlate}] {2$\times$ H2D, 1$\times$ D2H} (g3.west);

      \node[sl_good, text width=2.2cm] (v4) at (0, -4.55) {\textbf{4. Totalmente}\\\textbf{em GPU}};
      \node[sl_cpu, text width=3.5cm] (c4) at (3.6, -4.55) {Orquestração + sorteios\\($\approx 11{,}26$ KB/ger.)};
      \node[sl_good, text width=3.9cm] (g4) at (10.3, -4.55) {\textbf{Todo o ciclo residente}\\($D, N_k, P(t)$ + \texttt{RawKernel})};
      \draw[->, >=Stealth, line width=0.9pt, draw=ufscGreen] (c4.east) -- node[above=1pt, font=\tiny\bfseries\color{ufscGreen}] {Carga única + sementes} node[below=1pt, font=\tiny\color{ufscGreen}] {$0{,}00$ MB D2H (rotas)} (g4.west);
    \end{tikzpicture}%
  }{Elaborado pelo autor (taxonomia de \textcite{crainic2003parallel,crainic2010parallel}; híbridas idênticas bit a bit).}
\end{frame}

% --- Conjunto Experimental e Protocolo Estatístico --------------------------
\begin{frame}{Conjunto Experimental e Protocolo Estatístico}
  \begin{columns}[T,onlytextwidth]
    \begin{column}{0.46\textwidth}
      \begin{block}{\totalinstancias\ Instâncias da TSPLIB \autocite{reinelt1991tsplib}}
        \begin{itemize}
          \item \textbf{Dimensão:} \menorinstancia\ (\texttt{att48}) a \maiorinstancia\ (\texttt{pr1002}) cidades.
          \item \textbf{Distâncias:} \texttt{EUC\_2D} (10), \texttt{ATT} (2), \texttt{GEO} (2) e \texttt{EXPLICIT} (1).
          \item \textbf{Amostragem:} $\totalinstancias \times \totalvariantes \times \totalrepeticoes = \mathbf{\totalexecucoes}$ execuções (sementes pareadas \sementeinicial--\sementefinal).
        \end{itemize}
      \end{block}
    \end{column}
    \begin{column}{0.51\textwidth}
      \begin{block}{Protocolo de \textcite{demsar2006statistical} ($\alpha = 0{,}05$)}
        \textbf{Eixo 1 — Pareado por instância (\totalrepeticoes\ reps):}
        \begin{itemize}
          \item \textbf{Wilcoxon} nos 6 pares 2 a 2 \autocite{wilcoxon1945individual} + ajuste de \textbf{Holm--Bonferroni} ($90$ testes/métrica) \autocite{holm1979simple}.
        \end{itemize}
        \vspace{0.1cm}
        \textbf{Eixo 2 — Global multi-instância ($N=15, k=4$):}
        \begin{itemize}
          \item \textbf{Friedman} com correção $F_F$ \autocite{friedman1937use}; se $p_F < 0{,}05$, pós-teste de \textbf{Nemenyi} ($CD = 1{,}211$) \autocite{nemenyi1963distribution}.
        \end{itemize}
      \end{block}
    \end{column}
  \end{columns}
\end{frame}


% ============================================================================
% slides/04-resultados.tex — Capítulo 4: Resultados e Análise de Gargalos
% ============================================================================

\section{Resultados e Análise de Gargalos}

% --- Desempenho Agregado (15 Instâncias) ------------------------------------
\begin{frame}{Desempenho Agregado nas \totalinstancias\ Instâncias da TSPLIB}
  \begin{columns}[T,onlytextwidth]
    \begin{column}{0.32\textwidth}
      \kpicard{71,03$\times$}{Ganho Geom.\ vs.\ CPU}{Totalmente em GPU (mediana 0,34\,s)}
    \end{column}
    \begin{column}{0.32\textwidth}
      \kpicard{0,90\%--0,91\%}{\emph{Gap} Médio Preservado}{Sem dif.\ estatística ($p = 0{,}4543$)}
    \end{column}
    \begin{column}{0.32\textwidth}
      \kpicard{15\,/\,15}{Vitórias em Tempo}{Menor tempo em todas as instâncias}
    \end{column}
  \end{columns}

  \begin{slidetabela}[0.40\textheight]
    \begin{tabular}{@{}lccc@{}}
      \toprule
      \textbf{Algoritmo} & \textbf{Tempo mediano $[Q_{25};\, Q_{75}]$ (s)} & \textbf{\emph{Gap} médio (\%)} & \textbf{Ganho rel.\ CPU (geom.)} \\
      \midrule
      CPU (\emph{baseline})      & $23{,}62\;[4{,}81;\, 130{,}73]$        & $0{,}90 \pm 0{,}41$ & $1{,}00\times$ \\
      Híbrida Ingênua            & $4{,}81\;[1{,}17;\, 25{,}76]$          & $0{,}90 \pm 0{,}44$ & $4{,}57\times$ \\
      Híbrida Otimizada          & $1{,}02\;[0{,}18;\, 6{,}02]$           & $0{,}90 \pm 0{,}44$ & $23{,}86\times$ \\
      \rowcolor{ufscGreenLight}
      \textbf{Totalmente em GPU} & $\mathbf{0{,}34\;[0{,}05;\, 2{,}72]}$  & $0{,}91 \pm 0{,}46$ & $\mathbf{71{,}03\times}$ \\
      \bottomrule
    \end{tabular}
  \end{slidetabela}
  \notatabelaSlide{Elaborado pelo autor (\totalexecucoes\ execuções; pico de $112{,}48\times$ em \texttt{gr96}; em \texttt{pr1002}: $682{,}58\,\text{s} \to 14{,}03\,\text{s}$).}
\end{frame}

% --- Escalonamento em Tempo e Qualidade -------------------------------------
\begin{frame}{Escalonamento em Tempo de Execução e Qualidade (\emph{Gap})}
  \begin{columns}[c,onlytextwidth]
    \begin{column}{0.66\textwidth}
      \slideFigure[\slideTallFigH]{\linewidth}{scaling_time_gap.pdf}{Elaborado pelo autor (mediana e $Q_{25}\text{--}Q_{75}$ em \totalrepeticoes\ repetições).}
    \end{column}
    \begin{column}{0.32\textwidth}
      \begin{block}{Leitura das Curvas}
        \begin{itemize}
          \item \textbf{Tempo (log):} separação estável por ordens de grandeza até $n=\maiorinstancia$.
          \item \textbf{\emph{Gap}:} curvas praticamente sobrepostas entre as \totalvariantes\ variantes.
          \item \textbf{\texttt{berlin52}:} ótimo exato ($\text{gap}=0\%$) em 95/120 execuções.
          \item \textbf{\texttt{att532} $\to$ \texttt{si535}:} \texttt{si535} converge em $\approx 63$ ger.\ vs.\ 173--194 ger.\ em \texttt{att532}.
        \end{itemize}
      \end{block}
    \end{column}
  \end{columns}
\end{frame}

% --- Validação Estatística (Diagramas CD de Demšar) -------------------------
\begin{frame}{Validação Estatística: Tempo vs.\ Qualidade \autocite{demsar2006statistical}}
  \begin{columns}[T,onlytextwidth]
    \begin{column}{0.48\textwidth}
      \centering
      {\footnotesize\bfseries\color{ufscBlue} Tempo ($\chi_F^2 = 45{,}00$, $p < 0{,}001$)}
      \slideTikz[0.36\textheight]{\input{\chfourdir/tables/chapter4_validation_cd_diagram_time.tikz}}{}
      \begin{exampleblock}{Superioridade Temporal}
        \begin{itemize}
          \item Ordenação idêntica nas \totalinstancias\ instâncias ($1{,}00$ a $4{,}00$).
          \item \textbf{Eixo 1:} \textbf{90/90} pares significativos ($p_{\text{Holm}} \le 2{,}05\times 10^{-7}$).
        \end{itemize}
      \end{exampleblock}
    \end{column}
    \begin{column}{0.49\textwidth}
      \centering
      {\footnotesize\bfseries\color{ufscDark} \emph{Gap} ($F_F(3,42) = 0{,}8897$, $p = 0{,}4543$)}
      \slideTikz[0.36\textheight]{\input{\chfourdir/tables/chapter4_validation_cd_diagram_gap.tikz}}{}
      \begin{block}{Qualidade Equivalente}
        \begin{itemize}
          \item Friedman não rejeita $H_0$; nenhum par supera $CD = 1{,}211$.
          \item \textbf{Eixo 1:} \textbf{0/90} pares significativos ($p_{\text{Holm}} \ge 0{,}3418$).
        \end{itemize}
      \end{block}
    \end{column}
  \end{columns}
  \fonteSlide{Elaborado pelo autor (pós-teste de Nemenyi, $CD = 1{,}211$, $\alpha = 0{,}05$).}
\end{frame}

% --- Divisão do Tempo por Etapa e Lei de Amdahl -----------------------------
\begin{frame}{Divisão do Tempo por Etapa e Lei de Amdahl}
  \begin{slidetabeladensa}[0.48\textheight]
    \begin{tabular}{@{}llrcccc@{}}
      \toprule
      \textbf{Instância} & \textbf{Algoritmo} & \textbf{Tempo (ms)} & \textbf{Cruz.\ (\%)} & \textbf{BL $2$-opt (\%)} & \textbf{Sinc.\ (\%)} & \textbf{Demais (\%)} \\
      \midrule
      \multirow{4}{*}{\textbf{berlin52}}
        & CPU               & 1\,594,1   & 3,2           & \textbf{96,5} & 0,0           & 0,3 \\
        & Híbrida Ingênua   & 510,7      & 10,5          & \textbf{88,7} & 0,0           & 0,7 \\
        & Híbrida Otimizada & 78,8       & \textbf{79,0} & 13,8          & 0,2           & 6,6 \\
        & Totalmente em GPU & 21,2       & 8,7           & 8,7           & \textbf{36,0} & 45,0 \\
      \midrule
      \multirow{4}{*}{\textbf{pr264}}
        & CPU               & 23\,616,8  & 1,8           & \textbf{98,1} & 0,0           & 0,1 \\
        & Híbrida Ingênua   & 4\,813,3   & 14,4          & \textbf{84,6} & 0,0           & 1,0 \\
        & Híbrida Otimizada & 1\,020,2   & \textbf{60,9} & 35,3          & 0,1           & 3,5 \\
        & Totalmente em GPU & 337,7      & 2,5           & 2,6           & \textbf{84,7} & 9,9 \\
      \midrule
      \multirow{4}{*}{\textbf{pr1002}}
        & CPU               & 682\,575,6 & 1,4           & \textbf{98,4} & 0,0           & 0,2 \\
        & Híbrida Ingênua   & 115\,353,9 & 9,2           & \textbf{89,7} & 0,0           & 1,1 \\
        & Híbrida Otimizada & 25\,548,3  & \textbf{40,7} & \textbf{56,6} & 0,1           & 2,5 \\
        & Totalmente em GPU & 14\,027,6  & 0,9           & 0,7           & \textbf{94,8} & 3,4 \\
      \bottomrule
    \end{tabular}
  \end{slidetabeladensa}
  \notatabelaSlide{Elaborado pelo autor (médias em \totalrepeticoes\ repetições; demais = Inic., Sel., Mut., Aval., Sobr.).}

  \vspace{0.06cm}
  \begin{itemize}
    \footnotesize
    \item \textbf{Inversão na Híbrida Otimizada:} ao acelerar o $2$-opt em lote, o \textbf{cruzamento OX em CPU} sobe para \textbf{40,7\%--79,5\%} (média 62,2\%).
    \item \textbf{Papel de \texttt{Sinc.}:} como \emph{kernels} CUDA são assíncronos, \texttt{Sinc.} (36,0\%--94,8\%) mede a \textbf{execução efetiva na GPU}, não tempo ocioso.
  \end{itemize}
\end{frame}

% --- Tráfego PCIe e Lançamentos de Kernel -----------------------------------
\begin{frame}{Tráfego pelo Barramento PCIe e Lançamentos de \emph{Kernel}}
  \begin{slidetabeladensa}[0.46\textheight]
    \begin{tabular}{@{}llrrrr@{}}
      \toprule
      \textbf{Instância} & \textbf{Algoritmo} & \textbf{H2D (MB)} & \textbf{D2H (MB)} & \textbf{Total PCIe (MB)} & \textbf{Lançamentos} \\
      \midrule
      \multirow{3}{*}{\textbf{pr264}}
        & Híbrida Ingênua   & $22{,}54$    & $11{,}08$  & $33{,}62$    & $104\,960$ \\
        & Híbrida Otimizada & $23{,}30$    & $11{,}17$  & $34{,}47$    & $452$ \\
        & Totalmente em GPU & $1{,}55$     & $\mathbf{0{,}00}$ & $1{,}55$ & $532$ \\
      \midrule
      \multirow{3}{*}{\textbf{att532}}
        & Híbrida Ingênua   & $297{,}18$   & $96{,}53$  & $393{,}72$   & $453\,632$ \\
        & Híbrida Otimizada & $197{,}05$   & $96{,}90$  & $293{,}95$   & $1\,950$ \\
        & Totalmente em GPU & $6{,}17$     & $\mathbf{0{,}00}$ & $6{,}17$ & $2\,723$ \\
      \midrule
      \multirow{3}{*}{\textbf{pr1002}}
        & Híbrida Ingênua   & $2\,669{,}68$ & $543{,}26$ & $3\,212{,}94$ & $1\,355\,435$ \\
        & Híbrida Otimizada & $1\,099{,}67$ & $544{,}34$ & $1\,644{,}02$ & $5\,825$ \\
        \rowcolor{ufscGreenLight}
        & \textbf{Totalmente em GPU} & $18{,}69^*$ & $\mathbf{0{,}00}$ & $18{,}69^*$ & $6\,880$ \\
      \bottomrule
    \end{tabular}
  \end{slidetabeladensa}
  \notatabelaSlide{Elaborado pelo autor. $^*$Em \texttt{pr1002}, o tráfego físico real no PCIe é \textbf{10,66 MB} (8,03 MB contábeis do ponteiro de $D$).}

  \vspace{0.06cm}
  \begin{columns}[T,onlytextwidth]
    \begin{column}{0.48\textwidth}
      \begin{block}{Lote vs.\ Envio Individual}
        Híbrida Otimizada corta lançamentos em \textbf{$\sim 230\times$} ($1{,}35\text{M} \to 5{,}8\text{k}$ em \texttt{pr1002}) e evita reenvio $\mathcal{O}(n^2)$ de $D$.
      \end{block}
    \end{column}
    \begin{column}{0.49\textwidth}
      \begin{exampleblock}{Residência em VRAM}
        Totalmente em GPU zera retorno D2H (\textbf{0,00 MB}) e reduz o tráfego físico PCIe em \textbf{$301\times$} vs.\ Ingênua e \textbf{$154\times$} vs.\ Otimizada.
      \end{exampleblock}
    \end{column}
  \end{columns}
\end{frame}


% ============================================================================
% slides/05-discussao-conclusoes.tex — Capítulos 5 e 6: Discussão e Conclusões
% ============================================================================

\section{Discussão, Conclusões e Trabalhos Futuros}

% --- Discussão: Estado da Arte e Limitações ---------------------------------
\begin{frame}{Discussão: Estado da Arte e Limitações}
  \begin{columns}[T,onlytextwidth]
    \begin{column}{0.48\textwidth}
      \begin{block}{Frente ao Estado da Arte}
        \begin{itemize}
          \item Resolvedores especializados como \textbf{LKH} \autocite{helsgaun2000effective} e \textbf{AG+EAX} \autocite{nagata2013powerful} atingem frequentemente $0\%$ de \emph{gap}.
          \item Aqui, o AG+$2$-opt foi mantido controlado para \textbf{isolar o efeito da paralelização} ($\text{gap} \approx 0{,}90\%$).
          \item Para $n \gg 10^3$, o termo quadrático $4n^2$ de $D$ passa a pressionar a VRAM \autocite{schulz2013gpu,oneil2015rethinking}.
        \end{itemize}
      \end{block}
    \end{column}
    \begin{column}{0.49\textwidth}
      \begin{alertblock}{Ameaças à Validade}
        \begin{itemize}
          \item \textbf{Parada oracular:} $\text{gap} \le 1\%$ usa $B^*$ (\emph{time-to-target}); uso industrial depende de estagnação/tempo.
          \item \textbf{Desempate no $2$-opt:} reduções paralelas em GPU podem desempatar movimentos de igual $\Delta C$ distinto da CPU.
          \item \textbf{Telemetria e hardware único:} \texttt{Sinc.} agrega execução assíncrona; razões numéricas dependem do par i7-7700HQ + GTX 1050.
        \end{itemize}
      \end{alertblock}
    \end{column}
  \end{columns}
\end{frame}

% --- Implicações Práticas na Engenharia Mecânica ----------------------------
\begin{frame}{Implicações Práticas na Engenharia Mecânica e Pesquisa Operacional}
  \begin{block}{Aplicabilidade Direta}
    Estruturas de sequenciamento em \textbf{furação/corte CNC}, \textbf{robôs de solda} e \textbf{logística fabril} beneficiam-se de buscas aceleradas para avaliar mais cenários ou replanejar em tempo próximo do real.
  \end{block}

  \vspace{0.1cm}
  {\small\textbf{Cinco diretrizes práticas sintetizadas \autocite{kirk2010programming,schulz2013gpu,oneil2015rethinking,nvidia2024cuda}:}}
  \begin{enumerate}
    \small
    \item \textbf{Manter o estado na VRAM:} evitar envios individuais; usar acesso coalescido, memória compartilhada e mapas de bits.
    \item \textbf{Equilibrar o ciclo (Lei de Amdahl):} etapas leves em CPU (como OX) viram gargalo após acelerar o $2$-opt.
    \item \textbf{Reformular desvios para SIMT:} trocar ramificações irregulares por operações matriciais e reduções paralelas.
    \item \textbf{Explorar geometria:} listas de vizinhos ($k_{\text{viz}}=\kviz$) ampliam o alcance do hardware sem perda de qualidade.
    \item \textbf{Protocolo controlado:} aquecimento (\emph{warmup}) fora da medição, sementes fixas e testes não paramétricos.
  \end{enumerate}
\end{frame}

% --- Conclusões e Contribuições ---------------------------------------------
\begin{frame}{Conclusões e Contribuições}
  \begin{exampleblock}{Resposta à Questão de Pesquisa}
    Sob a mesma especificação algorítmica, a \textbf{estratégia de paralelização determina decisivamente o desempenho sem comprometer a qualidade}: a variante residente em VRAM atingiu \textbf{$71{,}03\times$} de ganho geométrico ($0{,}34\,\text{s}$ mediano) vs.\ $23{,}86\times$ e $4{,}57\times$ nas híbridas, com \emph{gaps} equivalentes ($0{,}90\%\text{--}0{,}91\%$; Friedman $p = 0{,}4543$).
  \end{exampleblock}

  \vspace{0.12cm}
  \begin{columns}[T,onlytextwidth]
    \begin{column}{0.48\textwidth}
      \begin{block}{Contribuições Empírica e Metodológica}
        \begin{itemize}
          \item Comparação controlada de \totalvariantes\ variantes sob o mesmo algoritmo memético.
          \item Protocolo reproduzível (\totalrepeticoes\ repetições pareadas e testes de Demšar).
        \end{itemize}
      \end{block}
    \end{column}
    \begin{column}{0.49\textwidth}
      \begin{block}{Contribuições Prática e Analítica}
        \begin{itemize}
          \item \emph{Framework} modular em três camadas com telemetria por etapa e PCIe.
          \item Caracterização do deslocamento de gargalos à luz da Lei de Amdahl.
        \end{itemize}
      \end{block}
    \end{column}
  \end{columns}
\end{frame}

% --- Trabalhos Futuros ------------------------------------------------------
\begin{frame}{Trabalhos Futuros}
  \begin{enumerate}
    \item \textbf{Operadores avançados em GPU:} adaptar busca local $3$-opt, movimentos tipo Lin--Kernighan e cruzamento EAX \autocite{nagata2013powerful} ao modelo SIMT para aproximar o \emph{gap} de $0\%$.
    \item \textbf{Instâncias de grande porte:} investigar dezenas de milhares de cidades com armazenamento triangular, recálculo sob demanda e particionamento de vizinhança \autocite{schulz2013gpu,oneil2015rethinking}.
    \item \textbf{Variantes industriais:} estender o \emph{framework} para TSP assimétrico (ATSP), roteamento capacitado (CVRP) e janelas de tempo.
    \item \textbf{Novas arquiteturas e energia:} avaliar CPUs \emph{multi-thread}, GPUs recentes e consumo energético.
    \item \textbf{Inferência bayesiana:} complementar testes frequentistas com testes bayesianos pareados \autocite{benavoli2017time}.
  \end{enumerate}
\end{frame}


% --- Referências (ABNT NBR 6023 compactas a partir de ../refs.bib) ----------
\begin{frame}[allowframebreaks]{Referências}
  \printbibliography[heading=none]
\end{frame}

% --- Encerramento / Perguntas da Banca --------------------------------------
\begin{frame}[standout]
  Obrigado pela atenção!\\[0.5cm]
  {\normalsize\normalfont Perguntas e considerações da Banca Examinadora}\\[0.7cm]
  {\footnotesize\normalfont
    Lucas Eduardo Galdino Silva \textbullet\ Orientador: Prof.~Dr.~Sérgio Fernando Mayerle\\
    Curso de Graduação em Engenharia Mecânica — CTC / UFSC
  }
\end{frame}

% --- Slides de apoio / Backup para arguição (fora da contagem principal) ----
\appendix
% ============================================================================
% slides/06-apendice-backup.tex — Slides de Apoio / Backup para Arguição
% (Excluídos da contagem principal via \appendix + appendixnumberbeamer)
% ============================================================================

\begin{frame}[standout]
  Apêndice — Slides de Apoio Técnico\\
  {\normalsize\normalfont Material suplementar para arguição da Banca Examinadora}
\end{frame}

% --- Backup 1: Especificações Completas de Software -------------------------
\begin{frame}{Backup 1: Especificações Completas do Ambiente de Software}
  \begin{slidetabela}
    \begin{tabular}{@{}lll@{}}
      \toprule
      \textbf{Componente} & \textbf{Versão / Requisito} & \textbf{Papel no ambiente} \\
      \midrule
      Sistema operacional            & Debian GNU/Linux 13, kernel 6.12.43 (\texttt{x86\_64}) & Plataforma do \emph{benchmark} \\
      \emph{Driver} NVIDIA           & 580.196.07                                             & Suporte ao CUDA Runtime 12.9 \\
      CUDA Toolkit / \texttt{nvcc}   & CUDA 12.9 (\texttt{nvcc} V12.9.86)                     & Compilação de \texttt{RawKernel} \\
      Compute Capability             & $\ge 5{,}0$ (testado: $6{,}1$, Pascal GP107)           & Requisito CUDA 12.9 / CuPy \\
      Compilador C/C++ (\emph{host}) & GCC 12.4.0                                             & Compilador auxiliar do \texttt{nvcc} \\
      Python                         & Python 3.12.13 ($\ge 3{,}12$)                          & Orquestração e interface \\
      Computação numérica            & NumPy 2.5.4; CuPy 14.2.0                               & \emph{Backends} CPU e GPU \\
      Estatística e dados            & SciPy 1.18.1; scikit-posthocs 0.17.1; DuckDB 2.0.0     & Protocolo Demšar e banco \\
      \bottomrule
    \end{tabular}
  \end{slidetabela}
  \notatabelaSlide{Elaborado pelo autor (reproduzível via \texttt{uv}, Seção~3.1.2 da monografia).}
\end{frame}

% --- Backup 2: Microbenchmark Lista de Candidatos (k_viz = 20) --------------
\begin{frame}{Backup 2: Microbenchmark da Lista de Candidatos ($k_{\text{viz}} = \kviz$)}
  \begin{slidetabela}[0.54\textheight]
    \begin{tabular}{@{}lrrrrrr@{}}
      \toprule
      \textbf{Instância} & \multicolumn{3}{c}{\textbf{CPU (s)}} & \multicolumn{3}{c}{\textbf{GPU — matriz de rotas (s)}} \\
      \cmidrule(lr){2-4} \cmidrule(lr){5-7}
       & $k_{\text{viz}} = \kviz$ & Exaustiva & Ganho & $k_{\text{viz}} = \kviz$ & Exaustiva & Ganho \\
      \midrule
      \texttt{berlin52} ($n=52$) & $0{,}935$ & $0{,}924$  & $0{,}99\times$         & $0{,}001$ & $0{,}001$ & $0{,}79\times$ \\
      \texttt{kroA100} ($n=100$) & $1{,}590$ & $4{,}250$  & $2{,}67\times$         & $0{,}002$ & $0{,}003$ & $1{,}51\times$ \\
      \texttt{ch150} ($n=150$)   & $3{,}122$ & $8{,}066$  & $2{,}58\times$         & $0{,}003$ & $0{,}005$ & $2{,}03\times$ \\
      \texttt{pr264} ($n=264$)   & $4{,}808$ & $39{,}513$ & \textbf{$8{,}22\times$} & $0{,}004$ & $0{,}016$ & $3{,}72\times$ \\
      \texttt{lin318} ($n=318$)  & $5{,}811$ & $45{,}798$ & $7{,}88\times$         & $0{,}006$ & $0{,}023$ & \textbf{$4{,}05\times$} \\
      \bottomrule
    \end{tabular}
  \end{slidetabela}
  \notatabelaSlide{Elaborado pelo autor (64 rotas, $\kappa = \maxpassadas$ varreduras; reduz complexidade de $\Theta(n^2)$ para $\Theta(n\,k_{\text{viz}})$).}

  \vspace{0.08cm}
  \begin{itemize}
    \footnotesize
    \item Ganho de até \textbf{$8{,}22\times$ na CPU} viabilizou executar integralmente a \emph{baseline} até $n = \maiorinstancia$ com \totalrepeticoes\ repetições.
  \end{itemize}
\end{frame}

% --- Backup 3: Resultados Detalhados nas 7 Instâncias Selecionadas ----------
\begin{frame}{Backup 3: Resultados Detalhados (7 Instâncias Selecionadas)}
  \begin{slidetabeladensa}[0.56\textheight]
    \begin{tabular}{@{}lcrrrrrr@{}}
      \toprule
      \textbf{Instância} & $n$ & \textbf{CPU (s)} & \textbf{Híb.\ Ing.\ (s)} & \textbf{Híb.\ Otim.\ (s)} & \textbf{Full GPU (s)} & \textbf{Ganho Full} & \textbf{\emph{Gap} Full (\%)} \\
      \midrule
      \texttt{berlin52} & 52   & $1{,}59$   & $0{,}51$   & $0{,}08$  & $\mathbf{0{,}02}$  & $75{,}03\times$ & $0{,}12$ \\
      \texttt{kroA100}  & 100  & $4{,}81$   & $1{,}18$   & $0{,}19$  & $\mathbf{0{,}05}$  & $94{,}86\times$ & $0{,}43$ \\
      \texttt{d198}     & 198  & $14{,}66$  & $3{,}08$   & $0{,}55$  & $\mathbf{0{,}17}$  & $88{,}62\times$ & $0{,}84$ \\
      \texttt{pr264}    & 264  & $23{,}62$  & $4{,}81$   & $1{,}02$  & $\mathbf{0{,}34}$  & $69{,}92\times$ & $0{,}80$ \\
      \texttt{att532}   & 532  & $158{,}24$ & $29{,}33$  & $6{,}67$  & $\mathbf{3{,}39}$  & $46{,}63\times$ & $1{,}13$ \\
      \texttt{si535}    & 535  & $84{,}19$  & $10{,}97$  & $2{,}48$  & $\mathbf{1{,}23}$  & $68{,}23\times$ & $0{,}95$ \\
      \texttt{pr1002}   & 1002 & $682{,}58$ & $115{,}35$ & $25{,}55$ & $\mathbf{14{,}03}$ & $48{,}66\times$ & $1{,}68$ \\
      \bottomrule
    \end{tabular}
  \end{slidetabeladensa}
  \notatabelaSlide{Elaborado pelo autor (médias de \totalrepeticoes\ repetições; ganhos calculados antes do arredondamento).}

  \vspace{0.08cm}
  \begin{itemize}
    \footnotesize
    \item Em \texttt{pr1002}: \emph{gaps} de $1{,}74\%$ (CPU), $1{,}64\%$ (híbridas, idênticas bit a bit) e $1{,}68\%$ (Totalmente em GPU); pico global em \texttt{gr96} ($112{,}48\times$).
  \end{itemize}
\end{frame}

% --- Backup 4: Eixo de Memória VRAM e Técnicas de Contorno ------------------
\begin{frame}{Backup 4: Consumo de VRAM e Contorno em Grande Porte}
  \begin{block}{Modelo de Consumo na Variante Totalmente em GPU}
    \[
      Mem_{\text{total}}(n) \approx \underbrace{4\,n^2}_{\text{matriz } D\;(\texttt{int32})} + \underbrace{2048\,n}_{2 \times 256 \times n \times 4\text{ B}} + \underbrace{2048}_{\text{custos }\texttt{int64}} \text{ bytes}
    \]
    \centering
    {\footnotesize Em \texttt{pr1002} ($n = \maiorinstancia$): $Mem_{\text{total}} \approx \mathbf{6{,}07\text{ MB}}$ ($\approx 0{,}15\%$ dos $4\text{ GB}$ de VRAM).}
  \end{block}

  \vspace{0.12cm}
  {\footnotesize \textbf{Estratégias quando $n \gg 10^3$ pressiona VRAM / memória compartilhada \autocite{schulz2013gpu,oneil2015rethinking}:}}
  \begin{columns}[T,onlytextwidth]
    \begin{column}{0.48\textwidth}
      \begin{itemize}
        \item \textbf{Armazenamento triangular:} explora simetria ($n(n-1)/2$, reduz memória pela metade).
        \item \textbf{Recálculo sob demanda:} troca VRAM por aritmética (penalidade temporal).
      \end{itemize}
    \end{column}
    \begin{column}{0.49\textwidth}
      \begin{itemize}
        \item \textbf{Particionamento da vizinhança:} fatiamento de movimentos $k$-opt.
        \item \textbf{Transferência em blocos / \emph{streams}:} sobrepõe cópias PCIe e computação.
      \end{itemize}
    \end{column}
  \end{columns}
\end{frame}

% --- Backup 5: Formulação Matemática do Protocolo de Demšar -----------------
\begin{frame}{Backup 5: Equações do Protocolo Estatístico \autocite{demsar2006statistical}}
  \begin{columns}[T,onlytextwidth]
    \begin{column}{0.48\textwidth}
      \begin{block}{Eixo 1: Wilcoxon + Holm--Bonferroni}
        Teste pareado de postos sinalizados \autocite{wilcoxon1945individual}:
        \[
          W = \min(R^+,\, R^-)
        \]
        Ajuste sequencial nos $m = 6$ pares ordenados $p_{(1)} \le \dots \le p_{(m)}$ \autocite{holm1979simple}:
        \[
          p_{(i)} < \frac{\alpha}{m - i + 1}, \quad \alpha = 0{,}05
        \]
      \end{block}
    \end{column}
    \begin{column}{0.49\textwidth}
      \begin{block}{Eixo 2: Friedman, Iman--Davenport e Nemenyi}
        Postos médios $R_j = \frac{1}{N}\sum_{i=1}^{N} r_i^j$ ($N = 15, k = 4$) \autocite{friedman1937use}:
        \[
          \chi_F^2 = \frac{12N}{k(k+1)}\!\left[\sum_{j=1}^k R_j^2 - \frac{k(k+1)^2}{4}\right], \quad
          F_F = \frac{(N-1)\chi_F^2}{N(k-1)-\chi_F^2}
        \]
        Diferença crítica de Nemenyi ($q_{0{,}05}=2{,}569$) \autocite{nemenyi1963distribution}:
        \[
          CD = q_{\alpha}\sqrt{\frac{k(k+1)}{6N}} = 1{,}211
        \]
      \end{block}
    \end{column}
  \end{columns}
\end{frame}



\end{document}
